\section{Zusammenfassung und Ausblick}
% [ ] Fassen sie die Beiträge ihre Visualisierungsanwendung zusammen. Wo bietet sie für die Personen der Zielgruppe einen echten Mehrwert.
% [ ] Was wären mögliche sinnvolle Erweiterungen, entweder auf der Ebene der Visualisierungen und/oder auf der Datenebene?

Das Ziel dieser Visual-Analytics-Anwendung war es, die Anwender beim kritischen Hinterfragen, der im Medaillenspiegel dargestellten Ergebnisse, zu unterstützen. Dazu sollten sie sich mit der Fragestellung auseinandersetzen, welches Land das sportlich erfolgreichste bei den Olympischen Sommerspielen 2024 war. Die Visualisierungstechniken helfen ihnen bei der Beantwortung, indem die erzielten Ergebnisse eines Landes aus einer neuen Perspektive darstellen werden. Dazu werden drei verschiedene Visualisierungen genutzt.

Mithilfe der ersten Visualisierung können die Anwender die Verteilung der Medaillen nach Sportarten und Disziplinen analysieren. Dies war mit dem bekannten Medaillenspiegel bisher nicht möglich. Dadurch können nun vielseitig erfolgreiche Länder genauso gut identifiziert werden, wie Länder, die den Großteil ihrer Medaillen nur in einer einzigen Sportart gewonnen haben. Die zweite Visualisierung präsentiert die Erfolge kleinerer Länder in Abhängigkeit von demografischen und wirtschaftlichen Einflussfaktoren. Diese alternativen Länderrankings werden mithilfe eines Parallelen Koordinatensystems vergleichbar gemacht. Den Anwendern könnte das Verhältnis \enquote{Medaillen pro Einwohnerzahl} als Bewertungsmaßstab zwar bekannt sein, aber eine Grafik, welche dieses Ranking im Vergleich zu dem bekannten Medaillenspiegel darstellt, wird in dieser Visual-Analytics-Anwendung neu eingeführt. Die dritte Visualisierung schafft es, den Anwendern Daten über die letzten 120 Jahre kompakt und übersichtlich aufzubereiten. Bereits diese Eigenschaft ist ein großer Mehrwert für Analysen rund um die Olympischen Sommerspiele. Zur Beantwortung der Fragestellung werden den Anwendern mithilfe von Farben Trends und Entwicklungen hervorgehoben, welche ohne die Visualisierung nicht erkennbar wären.

Aufgrund der verschiedenen Interaktionsmöglichkeiten können die Anwender die Visualisierungen nach ihren persönlichen Vorlieben verändern, um die, für sie am besten geeignete Ansicht zum Beantworten der Fragestellung, zu erhalten.

Die verwendeten Datensätze enthalten neben den Informationen über die Medaillen auch konkrete Daten von den erzielten Leistungen der Sportler. Dadurch könnte zum Beispiel die Abstände zwischen den Erstplatzierten untersucht werden. Die Anwendungsaufgaben könnten dann um die Frage nach dem sportlich erfolgreichsten Sportler erweitert werden. Andererseits ermöglicht die Visualisierung der Weltkarte in \cite{relWork2} völlig neue Analysemöglichkeiten. Dadurch könnte die Frage nach globalen Tendenzen beantwortet werden. 